% Standalone problem document, generated from the adjacent README.md.
% Compile with XeLaTeX. Mathematical content is maintained in Markdown.
\documentclass[11pt,a4paper]{article}
\usepackage[fontsize=10.5pt]{scrextend}
\usepackage[margin=22mm,headheight=15pt,headsep=9mm,footskip=11mm]{geometry}
\usepackage{fontspec}
\setmainfont{texgyrepagella}[Extension=.otf,UprightFont=*-regular,BoldFont=*-bold,ItalicFont=*-italic,BoldItalicFont=*-bolditalic]
\setsansfont{texgyreheros}[Extension=.otf,UprightFont=*-regular,BoldFont=*-bold,ItalicFont=*-italic,BoldItalicFont=*-bolditalic]
\setmonofont{texgyrecursor}[Extension=.otf,UprightFont=*-regular,BoldFont=*-bold,ItalicFont=*-italic,BoldItalicFont=*-bolditalic,Scale=MatchLowercase]
\usepackage{amsmath,amssymb}
\usepackage{unicode-math}
\setmathfont{texgyrepagella-math.otf}
\usepackage{xcolor,microtype,enumitem,needspace,fancyhdr}
\usepackage{bookmark}
\definecolor{ink}{HTML}{17344B}
\definecolor{accent}{HTML}{197D80}
\definecolor{muted}{HTML}{596773}
\hypersetup{colorlinks=true,linkcolor=accent,urlcolor=accent,pdftitle={MI-19 - A
q-permanent inequality for subset-preserving
permutations},pdfauthor={Open Problems in Numerical Linear Algebra}}
\urlstyle{same}
\setlength{\parindent}{0pt}
\setlength{\parskip}{5pt plus 1pt minus 1pt}
\linespread{1.04}
\setlength{\emergencystretch}{3em}
\widowpenalty=10000
\clubpenalty=10000
\displaywidowpenalty=10000
\allowdisplaybreaks[1]
\setlist{nosep,leftmargin=1.5em}
\providecommand{\tightlist}{\setlength{\itemsep}{0pt}\setlength{\parskip}{0pt}}
\providecommand{\pandocbounded}[1]{#1}
\pagestyle{fancy}
\fancyhf{}
\fancyhead[L]{\sffamily\footnotesize\color{muted}OPEN PROBLEMS IN NUMERICAL LINEAR ALGEBRA}
\fancyhead[R]{\sffamily\bfseries\footnotesize\color{accent}MI-19}
\fancyfoot[L]{\parbox[b]{0.9\textwidth}{\sffamily\fontsize{8}{10}\selectfont\color{muted}Editorial ratings; a bounded search cannot certify that a problem remains unsolved.\\Literature check: 2026-09-12}}
\fancyfoot[R]{\sffamily\footnotesize\color{muted}\thepage}
\renewcommand{\headrulewidth}{0.3pt}
\renewcommand{\headrule}{\hbox to\headwidth{\color{accent}\leaders\hrule height \headrulewidth\hfill}}
\makeatletter
\renewcommand{\section}{\Needspace{5\baselineskip}\@startsection{section}{1}{0pt}{12pt plus 2pt minus 2pt}{4pt}{\sffamily\bfseries\large\color{ink}}}
\renewcommand{\subsection}{\Needspace{4\baselineskip}\@startsection{subsection}{2}{0pt}{9pt plus 1pt minus 1pt}{3pt}{\sffamily\bfseries\normalsize\color{ink}}}
\makeatother
\setcounter{secnumdepth}{0}
\begin{document}
{\sffamily\small\color{muted}Matrix inequalities and norms\par}
\vspace{3pt}
{\raggedright\hyphenpenalty=10000\exhyphenpenalty=10000\sffamily\bfseries\fontsize{21}{24}\selectfont\color{ink}A
q-permanent inequality for subset-preserving permutations\par}
\vspace{7pt}
{\sffamily\small\textbf{Difficulty:} challenging\quad\textbf{Importance:} interesting
to specialist\par}
{\sffamily\small\bfseries\color{accent}Status: Lean verified\par}
\vspace{4pt}
{\color{accent}\hrule height 0.8pt}
\vspace{6pt}
\textbf{Rating rationale:} Arbitrary subsets introduce inversion-order
interactions absent for an initial segment; a solution would advance
specialized block inequalities for generalized permanents.

\subsection{Resolution --- 2026-09-11}\label{resolution-2026-09-11}

\textbf{Negative result by Matthew J. Colbrook}, Department of Applied
Mathematics and Theoretical Physics, University of Cambridge.
\textbf{Independent proof review: PASS.}

A real order-four PSD Gram matrix, \(q=7/8\) and the interior singleton
\(S=\{2\}\) give full minus restricted \(q\)-permanent equal to
\(-3235575/16384\). Inversions are counted in the full original
ordering. The strict counterexample also persists under sufficiently
small positive diagonal perturbations.

The exact target is resolved. The original statement and source evidence
are retained below; its former difficulty rating is historical.

\textbf{Lean verification --- 2026-09-12:} the complete negative target
passed independent statement and proof reviews, followed by Linux
Comparator and kernel checks.
\href{https://github.com/ajt60gaibb/OpenProblemsInNLA/blob/main/matrix-inequalities-and-norms/MI-19\#lean-proof-and-verification-evidence--2026-09-12}{Proof
and verification evidence}.

\textbf{Primary manuscript:}
\href{https://github.com/ajt60gaibb/OpenProblemsInNLA/blob/main/matrix-inequalities-and-norms/MI-19/solution.pdf}{complete
proof PDF},
\href{https://github.com/ajt60gaibb/OpenProblemsInNLA/blob/main/matrix-inequalities-and-norms/MI-19/solution.tex}{standalone
TeX}, Theorem 1.1 and its proof;
\href{https://github.com/ajt60gaibb/OpenProblemsInNLA/blob/main/matrix-inequalities-and-norms/MI-19/solution.md}{authorship
and scope}. The
\href{https://github.com/ajt60gaibb/OpenProblemsInNLA/blob/main/references/colbrook-matrix-2026-09-11/verification/reviews/MI-19-review.md}{independent
review} checks the full original argument and records its hash. The
draft was AI-assisted; that informal audit was independent agent
verification, not external human peer review or formal certification.
\href{https://github.com/ajt60gaibb/OpenProblemsInNLA/blob/main/references/colbrook-matrix-2026-09-11/README.md}{Submission
record}.

\subsection{Lean proof and verification evidence ---
2026-09-12}\label{lean-proof-and-verification-evidence-2026-09-12}

\textbf{Mathematical counterexample:} Matthew J. Colbrook, Department of
Applied Mathematics and Theoretical Physics, University of Cambridge.
\textbf{Lean formalization:} George Stepaniants, Department of Computing
and Mathematical Sciences, California Institute of Technology, Pasadena,
California, USA.

The
\href{https://github.com/sgstepaniants/OpenProblemsInNLA/blob/cd44ce9bcb84ebc79a1aa934918d1f76b2a9c6e7/matrix-inequalities-and-norms/MI-19/lean/NLA/MI19/Proof.lean}{Lean
proof} and
\href{https://github.com/sgstepaniants/OpenProblemsInNLA/blob/cd44ce9bcb84ebc79a1aa934918d1f76b2a9c6e7/matrix-inequalities-and-norms/MI-19/lean/Solution.lean}{public
declarations} are fixed at
\href{https://github.com/sgstepaniants/OpenProblemsInNLA/tree/cd44ce9bcb84ebc79a1aa934918d1f76b2a9c6e7/matrix-inequalities-and-norms/MI-19/lean}{proof
revision cd44ce9}. The selected declarations are:

\begin{itemize}
\tightlist
\item
  \textit{NLA.MI19.counterexample};
\item
  \textit{NLA.MI19.not\_subsetConjecture}.
\end{itemize}

The second theorem negates the complete original statement: every
\(n\ge2\), every complex Hermitian PSD matrix, every real \(q\in[0,1]\),
and every nonempty proper subset, with setwise preservation and
inversion counts in the full original ordering. All witness hypotheses
are proved. The actual complex PSD order-four Gram witness has \(q=7/8\)
and the interior singleton, giving full minus restricted q-permanent
\(-3235575/16384\). Exact rank, all-q polynomial identities, and the
positive-diagonal-perturbation extension are not among the formalized
claims.

The toolchain is Lean 4.33.1, with
\href{https://github.com/alerad/leancert/commit/621a43d7cf21f87872392a01e874f2f1dbddc926}{LeanCert
621a43d} and
\href{https://github.com/leanprover-community/mathlib4/commit/0df444a360eaa60ab8c11dca51a86af692955474}{mathlib
0df444a};
\href{https://github.com/sgstepaniants/OpenProblemsInNLA/blob/cd44ce9bcb84ebc79a1aa934918d1f76b2a9c6e7/matrix-inequalities-and-norms/MI-19/lean/lake-manifest.json}{all
dependencies are locked}. LeanCert checks one exact scalar inequality in
kernel mode. Both public theorems and the five audited internal
declarations have only \textit{propext}, \textit{Classical.choice}, and
\textit{Quot.sound} in their transitive axiom closure.

The catalog reviewed
\href{https://github.com/sgstepaniants/OpenProblemsInNLA/actions/runs/34703363593/job/103578942995}{successful
Linux verification on 12 September 2026}: fresh problem builds, real
sandbox isolation controls, exact Challenge/Solution comparison with no
definition holes, the standard-three axiom whitelist, and Lean
default-kernel replay. The pinned mathlib dependency cache was used.
Both requested theorems passed; sorry, custom-axiom, statement-mismatch,
invalid-kernel-proof, and native-trust rejection controls also executed
successfully. The
\href{https://github.com/ajt60gaibb/OpenProblemsInNLA/blob/main/matrix-inequalities-and-norms/MI-19/lean/verification/linux-2026-09-12/OPERATIONAL-REVIEW.md}{operational
audit},
\href{https://github.com/ajt60gaibb/OpenProblemsInNLA/blob/main/matrix-inequalities-and-norms/MI-19/lean/verification/linux-2026-09-12/artifacts/lean-MI-19/verify-20260912T155027Z-3941/comparator.log}{Comparator
log}, and
\href{https://github.com/ajt60gaibb/OpenProblemsInNLA/blob/main/matrix-inequalities-and-norms/MI-19/lean/verification/linux-2026-09-12/artifacts/lean-MI-19/verify-20260912T155027Z-3941/result.json}{result
receipt} retain the dated evidence and exact source hashes.

The
\href{https://github.com/ajt60gaibb/OpenProblemsInNLA/blob/main/matrix-inequalities-and-norms/MI-19/lean/README.md}{project
guide},
\href{https://github.com/ajt60gaibb/OpenProblemsInNLA/blob/main/matrix-inequalities-and-norms/MI-19/lean/formalization.yaml}{formalization
manifest}, and
\href{https://github.com/ajt60gaibb/OpenProblemsInNLA/blob/main/matrix-inequalities-and-norms/MI-19/lean/reviews}{independent
reviews} document attribution, fidelity, and scope. These are
independent AI-agent reviews, not external human peer review. On a Linux
host satisfying the
\href{https://github.com/ajt60gaibb/OpenProblemsInNLA/blob/main/tools/lean/HARNESS.md}{verification
prerequisites}, run the following from the repository root at the pinned
proof revision:

\begin{verbatim}
tools/lean/bootstrap.sh /tmp/nla-lean-tools
tools/lean/verify.sh \
  matrix-inequalities-and-norms/MI-19/lean /tmp/nla-lean-tools
\end{verbatim}

\subsection{Problem statement}\label{problem-statement}

Let \(n\ge2\), let \(A=(a_{ij})\in\mathbb C^{n\times n}\) be Hermitian
positive semidefinite, and let \(q\in[0,1]\). Define \[
\operatorname{inv}(\sigma)=\#\{(i,j):i<j,\ \sigma(i)>\sigma(j)\},\qquad
P_q(A)=\sum_{\sigma\in S_n}q^{\operatorname{inv}(\sigma)}
\prod_{i=1}^n a_{i,\sigma(i)},
\] with \(0^0=1\). Is it true that every nonempty proper subset
\(S\subset\{1,\ldots,n\}\) satisfies \[
P_q(A)\ge
\sum_{\substack{\sigma\in S_n\\ \sigma(S)=S}}
q^{\operatorname{inv}(\sigma)}\prod_{i=1}^n a_{i,\sigma(i)}?
\] Here \(\sigma(S)=S\) means setwise preservation. The inversion counts
on the right are taken in the full ordering \(1,\ldots,n\).

\subsection{Relevance and ratings}\label{relevance-and-ratings}

This is a block comparison for a matrix function interpolating
determinant and permanent. At \(q=1\) it reduces to a known permanental
block inequality. For general \(q\), ordering matters: the right side
must not be replaced by a product of q-permanents of the two principal
submatrices for an arbitrary subset.

\newpage
\subsection{References}\label{references}

\begin{itemize}
\tightlist
\item
  R. B. Bapat and A. K. Lal, \emph{Inequalities for the q-permanent},
  Linear Algebra and its Applications 197--198 (1994), 397--409
  (\href{https://doi.org/10.1016/0024-3795(94)90497-9}{original paper}).
\item
  C. M. da Fonseca, \emph{The \(\mu\)-permanent revisited} (2018), §5,
  Conjecture 3 and Theorem 5.1
  (\href{https://arxiv.org/pdf/1804.02231}{primary manuscript}).
\end{itemize}

\subsection{Status check --- 2026-09-10}\label{status-check-2026-09-10}

Da Fonseca explicitly states this conjecture, attributes it to Bapat and
Lal, and distinguishes proved initial-segment cases from arbitrary
subsets. Searches for ``q-permanent'', ``mu-permanent'', ``Lieb'',
``subset'', ``conjecture'', and ``proof'', including 2025--2026, found
no general resolution. The newest explicit formulation located is from
2018, so this entry has weaker recent status evidence than the other
candidates. The separately refuted permanent-on-top conjecture is not
included.

\textbf{Independent audit:} Independently rechecked da Fonseca
Conjecture 3 and Theorem 5.1 and searched for later subset-preserving
q-permanent results. Initial segments form a proved subfamily; arbitrary
subsets remain unresolved in the sources located.
\end{document}
